% S02: printed pages 41–43; master physical PDF130–132.
% TEXT-LEVEL CANDIDATE — NOT A PIXEL-VERIFIED RELEASE.
%
% Controlling source:
% 30_NALLINO_PARS_I_II_III_MASTER_1162P.pdf
% Expected SHA-256:
% 544c16b6355c9b74e281aded657d657224bff738366d5260e1a610d31b0d6297
%
% Body: NALLINO_TRANSLATION.
% Original numbered notes: NALLINO_APPARATUS.
% U-markers: separate modern checking register.
%
% Paragraph anchors are supplied here. Original line endings, italic
% spans, marginal locators, and exact mathematical typography still
% require comparison against the source images.
%
% Compile with LuaLaTeX.

\documentclass[10pt,a4paper]{article}

\usepackage[margin=20mm]{geometry}
\usepackage{fontspec}
\IfFontExistsTF{Libertinus Serif}
  {\setmainfont{Libertinus Serif}}
  {\IfFontExistsTF{FreeSerif}
    {\setmainfont{FreeSerif}}
    {\setmainfont{DejaVu Serif}}}

\usepackage{amsmath}
\usepackage{multicol}
\usepackage{hyperref}
\hypersetup{hidelinks}

\setlength{\parindent}{1em}
\setlength{\parskip}{3pt}

% Source-page heading; output pagination need not equal source pagination.
\newcommand{\SP}[3]{%
  \clearpage
  \hypertarget{AB01-PDF#1}{}%
  \begin{center}
    \textbf{#2\quad #3}
  \end{center}%
}

% Stable source-unit anchor.
\newcommand{\A}[1]{\hypertarget{#1}{}}

% Historical note call and historical note.
\newcommand{\NM}[1]{\textsuperscript{(#1)}}
\newcommand{\NN}[2]{%
  \par\noindent\textup{(#1)}\ #2%
}

% Modern uncertainty flag, distinguished from historical note numbers.
\newcommand{\U}[2]{%
  #2\textsuperscript{\texttt{[#1]}}%
}

% Explicit sexagesimal component encoding.
% The component readings and their printed order signs remain under U06.
\newcommand{\SX}[7]{%
  \ensuremath{%
    #1^\circ\,#2'\,#3''\,#4'''\,%
    #5^{\mathrm{iv}}\,#6^{\mathrm{v}}\,#7^{\mathrm{vi}}%
  }%
}

\begin{document}

% ================================================================
% MASTER PDF130 — PRINTED PAGE 41
% ================================================================

\SP{0130}{41}{AL-BATTANI OPUS ASTRONOMICUM}

\A{AB01-PDF0130-P01}
diei quae 365 diebus addebatur. Qua re de motu Solis non dubitavit,
ita ut Solem aliam sphaeram habere cuius centrum a centris duarum
sphaerarum differret, supponeret (b).

\A{AB01-PDF0130-P02}
Veteres plerumque haec ex observationibus aestivis, quae fiunt cum
Sol per punctum solstitiale aestivum transit, deprehenderunt; sed
non ita verae videntur ut observationes factae cum Sol per
alterutrum punctum aequinoctiale transit, et praesertim per punctum
aequinoctii autumni; hoc enim tempore aer clarior et purior est quam
tempore aequinoctii vernalis. Et quia Solis motus in declinatione,
quando per punctum solstitiale transit, tardus est, at quando per
puncta aequinoctialia iter facit, velocissimus, Ptolemaeus nonnisi
observationibus autumnis confisus est, et iuxta eas calculos
instituit. Una ex observationibus Hipparchi qua usus est et de
cuius veritate non dubitavit, ea est, qua, ut narrat\NM{1}, Solem
per punctum aequinoctii autumni transiisse comperit anno 178 ab
Alexandro mortuo, die tertia ex quinque diebus epagomenis, hora
mediae noctis cuius crastinum fuit quarta epagomenorum dies,
Alexandria in urbe\NM{2}. Post hoc Ptolemaeus, 285 annis Aegyptiis
transactis, iterum observavit, quam observationem ipse in libro suo
accuratissimam fuisse dicit; et Solis transitum per punctum
aequinoctii autumni comperit anno tertio Antonini regis, id est
anno 463\NM{3} ab Alexandro mortuo, die nona mensis Coptici Athyr,
una hora circiter post Solis ortum, Alexandria in urbe. Intervallum
inter duas observationes illas invenit esse veraciter 285 annorum
Aegyptiorum, 70 dierum, et $1/4+1/20\ [=3/10]$ diei, pro 71 diebus
$1/4$, qui ex quartis partibus integris in iis \U{U01}{287} annis
colligi debebant. Ratio igitur huius unius diei minus
\U{U02}{$1/20$} diei, --- quo tempus observationis praecessit tempus
quartae partis diei 365 dies superantis, ad 285 annos inter duas
illas observationes transactos, est ut ratio unius diei ad 300
annos; qua re tempus unius anni per has duas observationes
inventum, complectitur 365 dies $1/4$, detracto $1/300$ diei, quod
est una pars et quinta e 360 partibus\NM{4}.

\A{AB01-PDF0130-P03}
Narrat etiam Ptolemaeus, se observationes antiquas aestivas, quae
factae sint ante Hipparchum, accepisse, e quibus observationem esse
tempore Apseudis, regis [\U{U03}{ἄρχων}] Athenarum, factam, cum Sol
per punctum solstitiale aestivum transivisset anno 108\NM{5} ante
Alexandrum mortuum, matutino XXI diei mensis Coptici
Pharmouthi\NM{6}; deinde seipsum Solis transitum per punctum
solstitii aestivi observasse anno 463\NM{7} post Alexandrum
mortuum, XI die mensis Coptici Mesori, duabus horis circiter post
mediam noctem cuius crastinum dies XII fuit\NM{8}. Inter has duas
observationes intercesserunt fere 571 anni Aegyptii, 140 dies et
$1/2+1/3\ [=5/6]$ diei, pro 142 diebus $1/2+1/4\ [=3/4]$, qui ex
quartis partibus in iis annis collecti essent si in annis quartae
partes integrae fuissent. Invenit igitur tempus solstitii aestivi
tempus quartae integrae partis diei uno die et
$2/3+1/4\ [=11/12]$ praecessisse, quorum ratio ad 571 annos
memoratos est ut ratio duorum integrorum dierum ad 600\NM{9}
annos; idque cum eo iuxta quod calculos fecit convenit, scilicet
singulis 300 annis tempus observationis tempus quartae partis
integrae diei uno die praecedere, licet ob causam memoratam hae
observationes aestivae non adeo sint bonae ut autumnae. Manifestum
est inter primam illam observationem

\begin{multicols}{2}
\footnotesize

\A{AB01-PDF0130-N01}
\NN{1}{Almag. III, 2 (ed. Halma, t. I, p. 161).}

\A{AB01-PDF0130-N02}
\NN{2}{Al-Battāni coniectura; minime Ptolemaeus dicit Hipparchum
id Alexandriae observasse. Et re vera nunc certum habetur
Hipparchum Rhodi tantum observationes suas perfecisse, quam
perperam sub eodem meridiano quo Alexandria iacere putabat.}

\A{AB01-PDF0130-N03}
\NN{3}{Codex perperam 563.}

\A{AB01-PDF0130-N04}
\NN{4}{\texttt{[U04: AB01-PDF0130-N04-M01]}}

\A{AB01-PDF0130-N05}
\NN{5}{Codex perperam 168.}

\A{AB01-PDF0130-N06}
\NN{6}{Est observatio Metonis et Euctemonis
(27 Iun. 432 ante Chr. n.); Almag. III, 2
(ed. Halma, t. I, p. 162).}

\A{AB01-PDF0130-N07}
\NN{7}{Male Plato 36.}

\A{AB01-PDF0130-N08}
\NN{8}{Almag. l. l.}

\A{AB01-PDF0130-N09}
\NN{9}{Plato perperam 60.}

\end{multicols}

% ================================================================
% MASTER PDF131 — PRINTED PAGE 42
% ================================================================

\SP{0131}{42}{C. A. NALLINO}

\A{AB01-PDF0131-P01}
et observationem Hipparchi idem fere tempus intercessisse quod
inter Hipparchi et Ptolemaei; fuit enim observatio 286\NM{1} annis
ante Hipparchum.

\A{AB01-PDF0131-P02}
Nos ipsi in urbe ar-Raqqah observavimus; et una nostrarum
observationum autumnarum, qua confidimus et de cuius veritate,
ut ex instrumentis apparuit, haud dubitamus, est observatio quae
743 annis fuit post observationem autumnam Ptolemaei iam
memoratam. Ea invenimus Solem per punctum aequinoctii autumni
transivisse, anno 1194 aerae Dhū 'l-qarnayn, idest 1206 post
Alexandrum mortuum, ante Solis ortum diei XIX mensis Romani
Aylūl\NM{2}, scilicet VIII diei mensis Coptici Pachon, quatuor
circiter horis et $1/2+1/4$ (c). Et quia circulus meridianus
Alexandriae circulum meridianum urbis ar-Raqqah $2/3$ horae
aequinoctialis circiter praecedit\NM{3}, inter observationem nostram
et observationem autumnam Ptolemaei 743 anni Aegyptii, 178 dies,
et $1/2+1/4$ diei, detractis \U{U05}{$2/5$} horae circiter\NM{4},
intercesserunt, loco 185 dierum $1/2+1/4$ qui debuissent in his
annis colligi si quartae partes diei integrae fuissent. Si hos
7 dies et \U{U05}{$2/5$} horae, quibus tempus observationis tempus
quartae partis diei 365 dies superantis praecessit, per 743 annos
qui inter duas observationes illas fuerunt dividimus, invenimus
portionem unius anni esse $3^\circ24'$ ex iis $360^\circ$ qui
revolutionem diei et noctis perficiunt. Quam portionem si de
tempore quartae partis diei, scilicet de $90^\circ$ detrahimus,
quantitas augmenti super 365 dies integros remanet $86^\circ36'$;
est igitur verum anni tempus $365^\circ14'26''$\NM{5} fere.

\A{AB01-PDF0131-P03}
Si $360^\circ$ circumferentiae caelestis sphaerae per tempus anni
nuper inventum dividimus, motum medium Solis in nychthemero
reperimus \U{U06}{\SX{0}{59}{8}{20}{46}{56}{14}} esse, motum autem
in 30 diebus, qui mensem Aegyptium efficiant,
\U{U06}{\SX{29}{34}{10}{23}{28}{6}{47}}, et in 365 diebus anni
Aegyptii \U{U06}{\SX{359}{45}{46}{25}{32}{2}{31}}\NM{6} circiter
(d). Hos motus multiplicavimus et in tabulis ad annos collectos
et singulos, et ad menses, dies, horas, iuxta aeram Arabum
aeramque Romanorum descripsimus\NM{7}, ut facile quocumque harum
aerarum tempore locus Solis, secundum eius motum medium qui
medium Solis appellatur, inveniri possit. Tempus anni a nobis
inventum patet $2^\circ1/5$ minus esse quam tempus a Ptolemaeo
memorato\NM{8}; qua re motus Solis, iuxta calculum nostrum,
motum a Ptolemaeo inventum in die
\U{U06}{\SX{0}{0}{0}{3}{33}{43}{43}}\NM{9} superat, et in anno
Aegyptio, Deo volente,
\U{U06}{\SX{0}{0}{21}{40}{10}{53}{56}} fere\NM{10}.

\begin{multicols}{2}
\footnotesize

\A{AB01-PDF0131-N01}
\NN{1}{Codex 280, Plato 186.}

\A{AB01-PDF0131-N02}
\NN{2}{I. e. ante Solis ortum diei 19 Sept. 882; nam al-Battāni
prima die mensis Septembris annos huius aerae incipere vult,
cfr. cap. XXXII. Hunc locum laudat Ibn Yūnus apud Caussin,
\textit{Le livre de la grande table Hakémite}
(Notices et extraits des mss. de la Bibl. Nation., t. VII,
Paris 1804, p. 149).}

\A{AB01-PDF0131-N03}
\NN{3}{In suis tabulis geographicis, e communi Arabica traditione
depromptis, al-Battāni urbi ar-Raqqah $73^\circ15'$ long. tribuit,
Alexandriae, Ptolemaeum sequens, $60^\circ30'$. Hinc differentia
longitudinalis $12^\circ45'$ sequeretur, non $10^\circ$.}

\A{AB01-PDF0131-N04}
\NN{4}{I. e. 743 anni,
$178^{\mathrm d}\,17^{\mathrm h}\,36^{\mathrm m}$.}

\A{AB01-PDF0131-N05}
\NN{5}{I. e.
$365^{\mathrm d}\,5^{\mathrm h}\,46^{\mathrm m}\,24^{\mathrm s}$.
Ptolemaeus, Almag. III, 2 (Halma t. I, p. 165), seu potius
Hipparchus, invenerat
$365^\circ14'48''=
365^{\mathrm d}\,5^{\mathrm h}\,55^{\mathrm m}\,12^{\mathrm s}$;
astronomi nostri temporis
$365^{\mathrm d}\,5^{\mathrm h}\,48^{\mathrm m}\,47^{\mathrm s}{,}33$
ponunt.}

\A{AB01-PDF0131-N06}
\NN{6}{Male quartae in cod. et Plat. $31^{\mathrm{iv}}$.
Iam correxit Ed. Halley,
\textit{Emendationes ac notae in vetustas Albatênii observationes}
(Philosophical Transactions of the Royal Society, 1693,
vol. XVII, nr. 204, p. 916).}

\A{AB01-PDF0131-N07}
\NN{7}{Vide tabulas in fol. 164, r.--166, v.,
et 186, v.--189, r.}

\A{AB01-PDF0131-N08}
\NN{8}{Nam
$8^{\mathrm m}48^{\mathrm s}\times15
=528^{\mathrm s}\times15
=7920''
=2^\circ12'
=2^\circ1/5$.}

\A{AB01-PDF0131-N09}
\NN{9}{Codex male \U{U07}{$30^{[?]}$ et $42^{[?]}$};
Ptolemaeus III, 2 (ed. Halma, t. I, p. 166) enim habet
\U{U06}{\SX{0}{59}{8}{17}{13}{12}{31}}.}

\A{AB01-PDF0131-N10}
\NN{10}{Male quintae apud Platonem $55^{\mathrm v}$ et in codice
$50^{\mathrm v}$; Ptolemaeus
\U{U06}{\SX{359}{45}{24}{45}{21}{8}{35}}.}

\end{multicols}

% ================================================================
% MASTER PDF132 — PRINTED PAGE 43
% ================================================================

\SP{0132}{43}{AL-BATTANI OPUS ASTRONOMICUM}

\A{AB01-PDF0132-H01}
\begin{center}
\textbf{CAPUT XXVIII.}\par
De inaequalitate motus [sive anomalia] Solis et de eo quod
ab eius apogei loco cum illa apparet\NM{1}.
\end{center}

\A{AB01-PDF0132-P01}
Dicit [auctor]: Dissertatione de tempore anni et de motu medio
Solis absoluta, ad explicandum pergamus quae et quanta sit
inaequalitas in motu Solis, et qualis appareat [cum Sol motu suo
medio recesserit] a loco sui apogei in ecliptica.

\A{AB01-PDF0132-P02}
Rationem qua Ptolemaeus in libro suo\NM{2} usus est sequentes,
nos quo modo Sol quadrantes eclipticae percurrat, per plurimas
diligentissimas observationes multis subsequentibus annis factas
studuimus, invenimusque Solem eclipticam a puncto aequinoctii
autumni ad aequinoctii vernalis 178 diebus, 14 horis $1/2$
circiter percurrere; a puncto autem aequinoctii vernalis ad
aequinoctii autumni tempus eo longius occupare, quod ex nostris
diligentissimis observationibus patuit 186 dierum, 14 horarum
aequinoctialium et $1/4+1/2$ horae circiter esse. Et ex iis quae
memoravimus\NM{3} apparuit eius apogeum in hac dimidia parte haberi.

\A{AB01-PDF0132-P03}
Deinde observavimus et invenimus Solem spatium ab initio Arietis
ad initium Cancri, idest inter punctum aequinoctii vernalis et
punctum solstitii aestivi, 93 diebus et 14 fere horis
aequinoctialibus percurrere (a), quae parum ad deminutionem
inclinant; unde manifestum est cursui Solis a puncto aequinoctii
vernalis ad solstitii aestivi tempus longius necessarium esse
quam a puncto solstitii aestivi ad aequinoctii autumni.
Scimus igitur punctum apogei et centrum excentrici, super quem
punctum apogei simul ac super eclipticam incidit, ea quarta parte
contineri, quae tardioris temporis est quam reliquus quadrans.
Motum medium Solis in iis 186 diebus, 14 horis et $1/2+1/4$
horae invenimus $183^\circ56'12''$ esse, et in iis 93 diebus et
14 horis $92^\circ14'10''$ fere\NM{4}.

\A{AB01-PDF0132-P04}
Quae cum ita sint, circulum eclipticae ABCD circa centrum E
describamus. Duae diametri AC, BD se invicem ad angulos rectos
abscindant, sitque A punctum aequinoctii vernalis; erit B punctum
solstitii aestivi, C aequinoctii autumni, D solstitii hiemalis.
In quadrante AB, ob ea quae supra diximus, punctum F centrum
ponamus, et, intra primum circulum, circulum KLMN sphaerae
excentricae Solis describamus, cuius diametri KM, LN se invicem
in centro F ad angulos rectos abscindant. Punctum lineis BD, KM
commune littera S signemus; punctum quo diametrus AC circulum
KLMN versus A secat, littera P; punctum quo diametrus BD circulum
KLMN versus B abscindit, littera R. In arcu PK, perpendicularem
a P ad punctum Q diametri KM ducamus; item perpendicularem RG,
et lineam EF quae per duo centra transit, eamque usque ad
punctum H eclipticae ABCD protrahamus, littera T signantes locum
quo ea circulum KLMN secat.

\A{AB01-PDF0132-P05}
Manifestum est arcum AB esse $90^\circ$, ut arcum KL excentrici;
punctum P excentrici esse initium Arietis; R locum initii Cancri;
denique arcum \U{U08}{PKLRMY} excentrici eam excentrici

% The sentence continues on physical PDF133 / printed page 44.
% Do not supply a conjectural ending.

\begin{multicols}{2}
\footnotesize

\A{AB01-PDF0132-N01}
\NN{1}{I. e. de supputanda anomalia ad singulos gradus medii
motus Solis ab apogeo.}

\A{AB01-PDF0132-N02}
\NN{2}{Almag. III, 4 (ed. Halma t. I, p. 184--185.%
\textsuperscript{\texttt{[U09]}}}

\A{AB01-PDF0132-N03}
\NN{3}{Sed antea, ut me Schiaparelli monet, nihil de his rebus
dixit. Forte aliquid excidit ubi demonstrabatur apogeum
reperiendum esse in parte cui maius temporis intervallum
respondet.}

\A{AB01-PDF0132-N04}
\NN{4}{Ita etiam in calculis infra, qua re corrigere non audeo;
e tabulis $92^\circ14'26''$ deducuntur. Sed antea dixerat
«93 diebus et 14 horis aequinoctialibus, quae parum ad
deminutionem inclinant»; forte huic «parum minus» respondent
hae 16 secundae, quae in auctoris calculo deficere videntur.}

\end{multicols}

\end{document}